% ===================================================================
%  BẢNG Số 9 — Núi. Rừng.
%
%  Traduction, faite en 2026, en vietnamien standard d'aujourd'hui,
%  sur l'ido de texto/io. Regles de la colonne : voir 10-bang-01.tex.
%  Deux scenes, deux numerotations : les cles portent c1 et c2.
%
%  LE TABLEAU DE LA MONTAGNE EST CELUI OU LE SINO-VIETNAMIEN REPREND
%  LA MAIN, ET LA REGLE EST LA MEME QU'AU TABLEAU 1 : ce qui se
%  montre du doigt est vietnamien, ce qui s'apprend en classe est
%  chinois. Núi la montagne, đỉnh le sommet, hang la grotte, suối le
%  torrent, rừng la foret, cây le sapin, đá le rocher — tous
%  vietnamiens. Mais l'observatoire est « đài khí tượng », le
%  funiculaire « đường sắt leo núi », la station thermale « khu suối
%  khoáng », le sanatorium « nhà điều dưỡng », le haut fourneau « lò
%  cao » : ce sont des choses d'ingenieur, et l'ingenieur ecrivait en
%  caracteres.
%
%  ET LE GLACIER N'A PAS DE NOM, comme la neige du tableau 7 n'en
%  avait pas. « Sông băng » — le fleuve de glace — est un compose
%  transparent, calque sur le chinois, et c'est le seul mot dont
%  dispose une langue tropicale pour une chose qu'elle n'a jamais vue.
%  L'avalanche de meme : « tuyết lở », l'ecroulement de neige.
%
%  « BOHEMIEN » (41) A RECU AU QUEBEC UNE MAJUSCULE D'ETHNONYME ; ici
%  il en recoit une aussi, « người Rom », parce que le vietnamien n'a
%  jamais eu de mot familier pour ce peuple et n'a donc rien a
%  desapprendre. Le montreur d'ours et son ours ne bougent pas : ce
%  sont eux que la gravure montre.
%
%  ET LE PIC-VERT DU QUEBEC, QUI ETAIT LA UN PIC-BOIS, est ici « chim
%  gõ kiến », l'oiseau qui frappe — le nom dit ce qu'il fait, comme
%  partout ailleurs dans cette colonne.
% ===================================================================

%%K t09-apar-1 apar
\VUsaut{4.0mm}
\VUtitre{15.0pt}{60}{BẢNG Số 9}
\VUsaut{3.0mm}

%%K t09-tit-1 sub
\VUcentre{12.0pt}{0}{\textit{Cảnh thứ nhất.}}
\VUsaut{3.0mm}

%%K t09-tit-2 sub
\VUpk{11.4pt}{40}{Núi. --- Thị trấn suối khoáng.}
\VUsaut{3.0mm}

%%K t09-c1-01-1 p
1. --- Đã tám ngày em ở Pơrát. Em không sao nói hết lòng thán phục khi đứng trước
\VUgras{dãy núi} \textsuperscript{(1)} kỳ vĩ của rặng Pi-rê-nê, mà \VUgras{đường đỉnh}
\textsuperscript{(2)} vẽ trên nền trời xanh như một dải hoa văn khổng lồ.

%%K t09-c1-02-1 p
2. --- Em không ngờ những \VUgras{đỉnh núi} \textsuperscript{(3)} của rặng Pi-rê-nê lại
đạt \VUgras{độ cao} lớn đến thế, và em sững sờ trước những \VUgras{sông băng}
\textsuperscript{(5)} tuyệt đẹp cùng những \VUgras{chóp núi} \textsuperscript{(4)} phủ
tuyết, nơi những trận \VUgras{tuyết lở} khủng khiếp đổ xuống.

%%K t09-tit-3 sub
\VUsaut{3.0mm}
\VUpk{10.2pt}{40}{Phong cảnh núi non.}
\VUsaut{3.0mm}

%%K t09-c1-03-1 p
3. --- Ngay hôm sau khi đến, em đi thăm \VUgras{ngọn núi lửa} \textsuperscript{(6)} già
đã tắt ở gần Pơrát. Trên những mép \VUgras{miệng núi lửa} khô cằn và trơ trụi, người ta
còn thấy rất rõ dấu vết của dòng \VUgras{dung nham} đã chảy, nhiều thế kỷ trước, trên
\VUgras{sườn núi} \textsuperscript{(7)}. Em đã tìm được, trên một \VUgras{triền dốc},
mấy mảnh dung nham ấy.

%%K t09-c1-04-1 p
4. --- Pơrát nằm trên biên giới, và \VUgras{pháo đài} \textsuperscript{(8)} giữ
\VUgras{đèo} \textsuperscript{(27)} ngăn nước Pháp với Tây Ban Nha làm ta nhớ hẳn đến
những lâu đài cổ bên bờ sông Ranh, dường như đang rình \VUgras{mồi} từ trên đỉnh đá của
mình.

%%K t09-c1-05-1 p
5. --- Một con \VUgras{đại bàng} \textsuperscript{(9)} khổng lồ, có tổ không xa
\VUgras{đồn} \textsuperscript{(8)}, thường thu hút sự chú ý của em. Con \VUgras{chim săn
mồi} ấy, với chiếc \VUgras{mỏ} khỏe, thường mang về cho con nó một con cừu con hay dê
con mà nó quắp trong \VUgras{vuốt}. \VUgras{Cánh} nó lớn đến nỗi nó có thể lượn rất lâu
với gánh nặng của mình.

%%K t09-tit-4 sub
\VUsaut{3.0mm}
\VUpk{10.2pt}{40}{Người đi bộ đường dài.}
\VUsaut{3.0mm}

%%K t09-c1-06-1 p
6. --- Ở đây, chuyến dạo chơi thú vị nhất là chuyến đi đến \VUgras{thác nước}
\textsuperscript{(10)} Cau. \VUgras{Dòng nước đổ} rơi cuộn xoáy rồi biến thành một
\VUgras{con suối} đẹp trong \VUgras{rừng thông} \textsuperscript{(11)} ở phía tây thị
trấn. Chúng em leo qua khu rừng ấy đến tận \VUgras{đài khí tượng}
\textsuperscript{(12)}, và từ trên \VUgras{vọng lâu}, chúng em thấy cả \VUgras{toàn
cảnh} của vùng. \VUgras{Phong cảnh} thật kỳ diệu, và \VUgras{thiên nhiên} như đã ban cho
xứ này mọi kho báu mà nó có.

%%K t09-c1-07-1 p
7. --- Để xuống, chúng em đi \VUgras{đường sắt leo núi} \textsuperscript{(13)} và dừng ở
một \VUgras{ga} nhỏ, gần \VUgras{phế tích} \textsuperscript{(14)} của một \VUgras{lâu
đài phong kiến} xưa kia là nơi ở của các lãnh chúa Pơrát.

%%K t09-c1-08-1 p
8. --- Tòa lâu đài đáng thương! Nó nghĩ gì khi nhìn xuống dưới chân mình \VUgras{thị
trấn suối khoáng} \textsuperscript{(15)} xinh xắn, nay đã thành một \VUgras{khu nghỉ
dưỡng} thời thượng?

%%K t09-c1-08-2 p
\VUgras{Khí hậu} ở đó thật dễ chịu, \VUgras{nước khoáng} thật công hiệu, đến nỗi
\VUgras{việc chữa bệnh} ở \VUgras{nhà điều dưỡng} \textsuperscript{(17)} thành một niềm
vui thật sự.

%%K t09-tit-5 sub
\VUsaut{3.0mm}
\VUpk{10.2pt}{40}{Một thị trấn nghỉ dưỡng dễ chịu.}
\VUsaut{3.0mm}

%%K t09-c1-09-1 p
9. --- Vả lại, người ta đã làm mọi thứ để kỳ nghỉ được dễ chịu. Một chiếc \VUgras{cầu
cạn} \textsuperscript{(16)} lớn đã được xây để đặt đường sắt; \VUgras{công viên}
\textsuperscript{(18)} của \VUgras{khu suối khoáng}, nơi người ta tổ chức
\VUgras{hòa nhạc}, được bố trí thật duyên dáng. Nhiều \VUgras{biệt thự}
\textsuperscript{(19)} xinh đẹp đã mọc lên quanh sòng bạc \textsuperscript{(21)}, và một
\VUgras{con đường lát} \textsuperscript{(22)} được bảo trì rất tốt giúp việc đi lại dễ
dàng.

%%K t09-c1-10-1 p
10. --- Một \VUgras{bờ dốc} \textsuperscript{(23)} đơn giản ngăn \VUgras{con đường}
\textsuperscript{(20)} với chính \VUgras{thung lũng} \textsuperscript{(25)}, dưới đáy
thung lũng trải ra một \VUgras{cái hồ} \textsuperscript{(26)} nhỏ xinh xắn, nuôi một
\VUgras{con suối} \textsuperscript{(24)} trong vắt.

%%K t09-c1-11-1 p
11. --- Phía bên kia thung lũng, người ta khai thác một \VUgras{mỏ sắt}
\textsuperscript{(28)}. Em đã nhiều lần đến xem \VUgras{thợ mỏ} xuống \VUgras{giếng
mỏ}, và em còn vào cả \VUgras{đường hầm} nơi họ làm việc suốt ngày, đào lấy
\VUgras{quặng} chứa \VUgras{kim loại}.

%%K t09-c1-12-1 p
12. --- Một \VUgras{xưởng luyện kim} lớn đã được xây gần mỏ để dùng ngay những của cải
lấy lên từ đó. \VUgras{Lò cao} \textsuperscript{(29)}, đốt bằng \VUgras{than đá} vốn dồi
dào ở đây, cho phép làm ra \VUgras{gang} nhanh chóng và rẻ.

%%K t09-c1-13-1 p
13. --- Không xa mỏ, người ta đào một \VUgras{mỏ đá} \textsuperscript{(30)}. Không phải
\VUgras{đá hoa cương}, cũng không phải \VUgras{đá pooc-phia}, cũng chẳng phải \VUgras{sa
thạch} mà bác \VUgras{thợ đá} già đẽo ra bằng chiếc \VUgras{cuốc chim} của mình, mà chỉ
là \VUgras{đá xây} thường để làm nhà.

%%K t09-c1-14-1 p
14. --- Một trong những vẻ đẹp nhất của xứ này chắc chắn là cây \VUgras{thông}
\textsuperscript{(31)}. Khắp nơi, dưới đáy một \VUgras{khe núi} \textsuperscript{(32)},
bên bờ một \VUgras{dòng thác} \textsuperscript{(33)}, một \VUgras{vực sâu}
\textsuperscript{(34)} hay một bờ đá dựng, những chiếc lá kim mảnh của chúng điểm một
sắc xanh.

%%K t09-tit-6 sub
\VUsaut{3.0mm}
\VUpk{10.2pt}{40}{Đi bộ.}
\VUsaut{3.0mm}

%%K t09-c1-15-1 p
15. --- Em tận dụng kỳ nghỉ để đi \VUgras{dạo xa}. Những \VUgras{con đường mòn}
\textsuperscript{(35)} uốn lượn trên núi cho phép đi, không cần để ý đến quãng đường,
nhiều \VUgras{ki-lô-mét} và thường là nhiều chục ki-lô-mét.

%%K t09-c1-15-2 p
Những \VUgras{cột mốc} \textsuperscript{(36)} và \VUgras{cột chỉ đường}
\textsuperscript{(37)} đánh dấu các con đường, nơi ta thường gặp một chàng \VUgras{trai
miền núi} \textsuperscript{(38)} với chiếc \VUgras{bị} \textsuperscript{(40)} bên hông,
mang một con \VUgras{sóc đất} \textsuperscript{(39)}, hoặc một \VUgras{người Rom}
\textsuperscript{(41)} nào đó mà chiếc \VUgras{mũ lông} \textsuperscript{(42)} tương
phản kỳ lạ với chiếc \VUgras{áo trùm đầu} \textsuperscript{(43)} cũ rách. Anh ta dắt
bằng \VUgras{xích} một con \VUgras{gấu} \textsuperscript{(44)} đã thuần.

%%K t09-tit-7 sub
\VUpk{10.2pt}{40}{Leo núi.}
\VUsaut{3.0mm}

%%K t09-c1-16-1 p
16. --- Hầu như ngày nào em cũng đi với một \VUgras{người dẫn đường}
\textsuperscript{(48)} để \VUgras{leo núi}, và em rất tự hào khi, \VUgras{ba lô}
\textsuperscript{(47)} trên lưng và cây « \VUgras{gậy leo núi} » trong tay, em, như một
\VUgras{người leo núi} \textsuperscript{(46)} thực thụ, trèo lên những \VUgras{tảng đá}
\textsuperscript{(45)}.

%%K t09-c1-17-1 p
17. --- Từ trên đỉnh núi, người dưới đồng bằng trông không to hơn con kiến; một
\VUgras{đàn cừu} \textsuperscript{(50)} đang gặm cỏ trên một \VUgras{bãi chăn}
\textsuperscript{(49)} trông như một món đồ chơi trẻ con. Một cây \VUgras{thông}
\textsuperscript{(51)}, hay cả một \VUgras{ngôi nhà gỗ} \textsuperscript{(52)}, chỉ như
một vết nhỏ trên nền xanh.

%%K t09-c1-18-1 p
18. --- Trong những chuyến đi ấy, chúng em thường gặp trên \VUgras{lối mòn}
\textsuperscript{(53)} một \VUgras{người săn sơn dương} \textsuperscript{(54)} vác trên
vai con vật vừa bắn được.

%%K t09-c1-19-1 p
19. --- Em đã ép vào tập cây của mình một nhánh \VUgras{dương xỉ}
\textsuperscript{(55)} rất lạ và một cành \VUgras{thạch nam} \textsuperscript{(57)} hồng
đẹp mà em hái ở cửa một cái \VUgras{hang} \textsuperscript{(56)} nhỏ trong chuyến dạo
chơi vừa rồi.

%%K t09-tit-8 sub
\VUsaut{3.0mm}
\VUcentre{12.0pt}{0}{\textit{Cảnh thứ hai.}}
\VUsaut{3.0mm}

%%K t09-tit-9 sub
\VUpk{11.4pt}{40}{Rừng. --- Đi săn.}
\VUsaut{3.0mm}

%%K t09-c2-01-1 p
1. --- Hôm qua em đã trải một ngày tuyệt vời trong rừng. Sự nhộn nhịp ngự trị giữa các
loài vật và \VUgras{côn trùng} \textsuperscript{(5)} sống ở đó làm em rất thích thú.

%%K t09-c2-01-2 p
Con \VUgras{nhện} \textsuperscript{(1)} giăng \VUgras{mạng} \textsuperscript{(2)} giữa
hai cành để bắt con \VUgras{ruồi} \textsuperscript{(3)} bay qua, con \VUgras{ốc sên}
\textsuperscript{(4)} nhọc nhằn mang cái vỏ của mình, con bọ sừng hươu đang xem xét mồi,
con \VUgras{kiến} chăm chỉ tất bật bên \VUgras{tổ kiến} \textsuperscript{(6)} — tất cả
những sinh vật bé nhỏ ấy đi, đến, làm việc với một sự siêng năng lạ thường.

%%K t09-c2-02-1 p
2. --- Một con \VUgras{rắn} \textsuperscript{(7)} mà thoạt đầu em tưởng là \VUgras{rắn
độc}, nhưng chỉ là một con \VUgras{rắn nước} thường, nằm nấp dưới một gốc cây và thỉnh
thoảng thò cái đầu nhỏ thon ra, lúc nào cũng như sắp cắn. Trước mặt em, một con
\VUgras{sên trần} \textsuperscript{(8)} to nặng nề bò đi sau khi đã ăn hết một trong
những quả \VUgras{dâu rừng} \textsuperscript{(11)} thơm ngon mọc đầy ở đó.

%%K t09-c2-03-1 p
3. --- Mặt đất phủ đầy \VUgras{hoa rừng} : những cây \VUgras{anh thảo}
\textsuperscript{(9)}, những cây \VUgras{dừa cạn} \textsuperscript{(10)} và nhiều cây
khác em không biết tên nở giữa những \VUgras{búi cỏ} \textsuperscript{(12)}.

%%K t09-c2-04-1 p
4. --- Ở vùng này, \VUgras{nấm cục} gần như cũng thường như \VUgras{nấm}
\textsuperscript{(14)}, và một con \VUgras{lợn con} \textsuperscript{(13)} của một người
gác rừng đang tham lam tìm xem có kiếm được vài củ không.

%%K t09-tit-10 sub
\VUsaut{3.0mm}
\VUpk{10.2pt}{40}{Săn trộm.}
\VUsaut{3.0mm}

%%K t09-c2-05-1 p
5. --- Em đã chứng kiến \VUgras{đoạn kết} của một cảnh làm em rất thích thú. Nhờ cái mũi
thính của con \VUgras{chó săn} \textsuperscript{(15)}, \VUgras{người gác rừng}
\textsuperscript{(16)} vừa bắt quả tang một \VUgras{kẻ săn trộm} \textsuperscript{(22)}
đúng lúc hắn sắp mang đi \VUgras{con mồi} \textsuperscript{(17)} vừa bắn được. Thà bỏ mồi
còn hơn để bị bắt, kẻ săn trộm đã chạy trốn, và con \VUgras{thỏ rừng}
\textsuperscript{(18)}, con \VUgras{nai cái} \textsuperscript{(19)}, con \VUgras{hoẵng}
\textsuperscript{(20)} cùng con \VUgras{lợn rừng} \textsuperscript{(21)} nằm lại trên cỏ,
làm người gác rừng mừng rỡ.

%%K t09-c2-06-1 p
6. --- Sự đa dạng của cây cối trong rừng thật đáng kinh ngạc. Cây \VUgras{sồi dẻ}
\textsuperscript{(23)} và cây \VUgras{bạch dương} \textsuperscript{(26)}, cây
\VUgras{thông} cho \VUgras{nhựa} \textsuperscript{(27)}, cây \VUgras{sồi}
\textsuperscript{(29)} và nhiều cây khác nữa đan cành vào nhau.

%%K t09-c2-06-2 p
Dây \VUgras{thường xuân} \textsuperscript{(24)} bám vào thân những cây cao làm cho
\VUgras{cánh rừng già} \textsuperscript{(25)} một màu xanh óng ánh, hòa hợp dễ chịu với
sắc nhạt và dịu hơn của \VUgras{rêu} \textsuperscript{(31)}, nơi con \VUgras{chim gõ
kiến} \textsuperscript{(30)} tìm côn trùng.

%%K t09-tit-11 sub
\VUsaut{3.0mm}
\VUpk{10.2pt}{40}{Phong cảnh rừng.}
\VUsaut{3.0mm}

%%K t09-c2-07-1 p
7. --- Những cây sồi già làm em say mê. Bộ \VUgras{rễ} \textsuperscript{(32)} to xoắn
xuýt của chúng, nửa trồi lên khỏi mặt đất, cho chúng một vẻ mạnh mẽ và cường tráng tuyệt
vời, và em tự hỏi làm sao \VUgras{người chủ} \textsuperscript{(35)} lại đành cho đốn
chúng và giao chúng cho chiếc \VUgras{cưa} \textsuperscript{(34)} của \VUgras{người tiều
phu} \textsuperscript{(33)}. Những \VUgras{thân cây} \textsuperscript{(36)} đáng thương,
bị lột \VUgras{vỏ} \textsuperscript{(37)} hoặc bị chẻ ra bằng \VUgras{rìu bổ củi}
\textsuperscript{(38)} của \VUgras{người làm công nhật} \textsuperscript{(39)}, làm em
buồn lòng.

%%K t09-c2-08-1 p
8. --- Ngay gần đó, một \VUgras{người đốt than} \textsuperscript{(41)} già đang dùng
\VUgras{xẻng} của mình xếp một \VUgras{đống} \textsuperscript{(42)} than, và một bà cụ
mang đi một \VUgras{bó} \textsuperscript{(40)} \VUgras{củi khô} nhặt từ cành những cây
đã đốn.

%%K t09-tit-12 sub
\VUsaut{3.0mm}
\VUpk{10.2pt}{40}{Cuộc săn.}
\VUsaut{3.0mm}

%%K t09-c2-09-1 p
9. --- Em muốn vào nghỉ một lát ở \VUgras{nhà người gác rừng} \textsuperscript{(46)},
nhưng khi đến gần cái \VUgras{ao} \textsuperscript{(43)}, nơi một con \VUgras{thiên nga}
\textsuperscript{(45)} đẹp đang bơi, em nhận ra rằng một trận \VUgras{lụt}
\textsuperscript{(44)} mới đây đã biến con đường thành một \VUgras{đầm lầy} thực sự. Em
phải quay lại và, khi đi ngang cây \VUgras{tuyết tùng} \textsuperscript{(49)}, những cây
\VUgras{dương} \textsuperscript{(48)} và cây \VUgras{du} \textsuperscript{(47)} ở bìa
rừng, em thấy từ một \VUgras{lùm cây} \textsuperscript{(54)} lao ra một con \VUgras{hươu}
\textsuperscript{(50)} tuyệt đẹp, bị một \VUgras{bầy chó săn} \textsuperscript{(51)} hung
hãn đuổi theo, mà tiếng \VUgras{tù và} \textsuperscript{(53)} của những \VUgras{thợ săn
cưỡi ngựa} \textsuperscript{(52)} còn thúc thêm. Con vật đáng thương đã kiệt sức. Nó ngã
xuống chết cách em vài bước.

%%K t09-c2-10-1 p
10. --- Con trai người gác rừng, bị tiếng động của \VUgras{cuộc săn đuổi} lôi cuốn, ra
khỏi nhà và hai chúng em cùng trở vào rừng. Cậu đã thấy một con \VUgras{chim ác là}
\textsuperscript{(56)} trên cành một cây sồi già. Cậu nghĩ tổ của nó hẳn giấu trong
\VUgras{tán lá} \textsuperscript{(58)} và cậu muốn lấy cái tổ ấy.

%%K t09-c2-10-2 p
Nhanh nhẹn như một con \VUgras{sóc} \textsuperscript{(61)}, cậu trèo từ \VUgras{cành}
\textsuperscript{(55)} này sang cành khác, nhưng chẳng thấy gì và chỉ làm bay lên một
con \VUgras{quạ khoang} \textsuperscript{(60)} già đậu trên một \VUgras{nhánh}
\textsuperscript{(57)}.
\VUsaut{4.0mm}
\VUfilet{22mm}
